\section{Line follower and MicroMouse logic}
\label{sec:paths}
\textit{Builds on: grids, curves (\ref{sec:curved}), ray casting for line of sight (\ref{sec:rt}).}

The intruder rovers are built like two classic robot-contest machines. \textbf{Planning (MicroMouse):} the
40\,m$\times$40\,m area is divided into $80\times80$ cells of 0.5\,m; every cell hidden from the UAVs by an occluder
is a goal; a breadth-first flood fill spreads distances out from the goals, and the rover walks ``downhill'' to the
neighbour with the smaller value, as a MicroMouse descends its maze towards the centre. \textbf{Driving (line-follower
robot):} the planned cell path is smoothed into a virtual guide line; five reflectance sensors in front of the rover
report whether they are over the line, and the weighted sensor error drives a PD controller that sets the left and
right wheel speeds of a differential drive. The soldiers use A* on a $72\times72$ grid of 1\,m cells with 0.40\,m
clearance, smoothed into centripetal Catmull--Rom splines.
\begin{equation}
d(c)=1+\min_{n\in N_4(c)}d(n),\ d(\text{goal})=0;\qquad
e=\frac{\sum_{i=1}^{5}w_i s_i}{\sum_i s_i},\ \ u=K_p e+K_d\dot e,\ \ v_{L,R}=v\mp u
\label{eq:paths}
\end{equation}
with $w_i$ the sensors' lateral offsets, $s_i\in\{0,1\}$ their readings, $K_p=3.2$, $K_d=0.25$.

\begin{figure}[H]
\centering
\begin{subfigure}{0.49\linewidth}\raw{v5/plan_routes}{\linewidth}\caption{\key{M}: flood field, routes and guide lines.}\end{subfigure}\hfill
\begin{subfigure}{0.49\linewidth}\raw{scene/day_patrol_spline}{\linewidth}\caption{A soldier's A* patrol spline.}\end{subfigure}
\caption{Path debug views drawn by the program. Notice in (a) that the routes bend towards cover (cells the UAVs cannot
see) rather than taking the straight line, and in (b) that the spline rounds the grid path's corners while keeping
clear of the walls.}
\label{fig:paths}
\end{figure}

\begin{table}[H]
\centering\small
\begin{tabularx}{\linewidth}{@{}lXXl@{}}
\toprule
Method & Result & Cost & Used by\\\midrule
Flood fill (BFS) & distance to the nearest goal from every cell; any start descends greedily & $O(\text{cells})$ per replan, then $O(1)$ per step & rovers\\
A* & one shortest path, start to goal & expands far fewer cells with a good heuristic & soldiers\\
PD line following & smooth tracking of a drawn line with sensor feedback & 5 samples per tick & rovers\\
Catmull--Rom (centripetal) & smooth curve through the waypoints, no loops or cusps & a few evaluations per sample & soldiers\\
\bottomrule
\end{tabularx}
\caption{Differences between the path methods.}
\end{table}
\where{src/ai/HideController.{h,cpp} (flood fill, guide line, PD gains kKp / kKd), src/scene/Garrison.cpp SitePlanner (A*), src/math/Spline.cpp}
\critique{The rover's sensors are simulated against a line the planner drew, not against a painted line in the scene,
so the line-follower layer adds realistic lag and overshoot rather than new information. On top of it, a deliberate
lateral weave, speed wander and short hesitations (\code{weave}, \code{speedWander}, \code{hesitate}) keep the UAVs'
tracker from extrapolating a rover along one straight, constant-speed line.}
